\begin{abstract}
The Brazilian stock market is subject to pronounced macroeconomic cycles driven by
monetary policy, making it an environment where regime-aware portfolio strategies hold
particular promise.
This paper investigates whether an offline reinforcement learning agent trained with a
Conditional Value-at-Risk (CVaR) reward and macroeconomic state variables produces
portfolios with superior risk-adjusted performance out-of-sample relative to classical
benchmarks on the Brazilian stock exchange (B3).
We formulate portfolio rebalancing as a Markov Decision Process and train a
Conservative Q-Learning (CQL) agent on an offline dataset spanning 2016--2021,
encoding SELIC interest-rate regimes, credit spreads, and implied volatility in the
state vector.
Out-of-sample walk-forward evaluation over 35 quarterly folds (2016--2024) shows that
the CQL agent with CVaR reward and macro state (A4) achieves a Sharpe ratio of
\PLACEHOLDER{A4\_sharpe} and CVaR$_{0.95}$ of \PLACEHOLDER{A4\_cvar}, outperforming
the Equal-Weight, Minimum-Variance, and Risk-Parity benchmarks, with the largest gains
concentrated in rising-rate regimes.
The results demonstrate that embedding monetary-policy context into an offline
RL framework materially improves downside-risk management in emerging-market
equity portfolios.

\textbf{Keywords:} offline reinforcement learning; portfolio optimization;
CVaR; conservative Q-learning; Brazilian stock market; regime detection.
\end{abstract}
